Habiendo ya enumerado los métodos clásicos y más actuales del campo, y teniendo en cuenta su alto desempeño en esta tarea, es difícil con los conocimientos de principiante en el campo dar unos resultados mejores. Este TFG no busca superar modelos del estado del arte como \textit{Spleeter},\textit{Open-Unmix} o \textit{Demucs} si no que se limita a dar un acercamiento incial "sencillo" para aprender sobre la construcción de los sistemas que dan solución a este problema y dar un estudio experimental controlado dentro del alcance al estado del arte sobre diferentes funciones de pérdida.

En concreto se ha enfocado en lo siguiente: implementar una proceso completo, controlado y reproducible. Trabajar con STFTs y máscaras. Comparar el desempeño de diferentes funciones de pérdida. Evaluar dos y cuatro fuentes, implementar una interfaz de uso fácil y entendible para el usuario y analizar problemas reales de entrenamiento, inferencia y evaluación. En resumen, no se busca superar el estado del arte, si no simplemente proporcionar una comparativa de diferentes funciones de pérdida y construir un proceso completo.